% !TeX root = main.tex
\section{应用场景说明}

\subsection{目标用户与使用入口}

核心用户是持续更新中长篇小说的个人作者和需要多人审校的小型内容团队。用户通过创作会话提出目标，查看人物、大纲、正文和审稿结果，并在关键写入前确认变更。迁移首期以 JiuwenSwarm 会话与运行过程展示完成闭环，后续逐步对接正文编辑和差异查看等工作台能力。

\noindent\begin{tabularx}{\textwidth}{@{}p{0.17\textwidth}YY@{}}
\toprule
目标用户 & 具体痛点 & 系统提供的支持 \\
\midrule
长篇连载作者 & 更新压力大，人物状态和伏笔跨章维护困难 & 上下文装配、章节初稿、连续性核查、伏笔提示 \\
旧稿接续作者 & 资料散落，修改后难判断哪些事实仍然有效 & 旧稿结构确认、事实接纳、作品状态重建、版本核对 \\
小型编辑团队 & 审稿尺度不一，意见难定位，改稿结果难追踪 & 稿件评审、修订提案、版本关联与复评记录 \\
\bottomrule
\end{tabularx}

\subsection{场景一：设定复杂的连载小说持续更新}

\textbf{使用情境。}作者已有人物设定和若干章节，准备续写一个悬疑单元，希望既推进线索，又避免提前泄露真相。例如，第八章人物尚未得知钥匙去向，第九章草稿却写出了只有知情者才能采取的行动。

\textbf{Agent 协作。}Planner-Agent 根据近期创作目标确认本章所透露的线索和伏笔；Writer-Agent 读取包含前章正文、人物认知、世界观和每章大纲的上下文，形成初稿；Fact-Review-Agent 检查人物知识与事件因果，Plot-Review-Agent 核查线索的兑现与场景节奏。涉及矛盾的意见须引用本章原文和对应事实依据，交由 Revision-Agent 形成改稿提案。

\textbf{作者获得的结果。}一份章节候选稿、一份审稿报告，以及可逐项判断的修订差异。作者可以保留自己有意安排的误导，也可以接受修复；接受后重新核查受影响内容，并同步更新作品状态。

\textbf{价值与验证。}重点观察人工补充设定的次数、连续性错误检出情况、有效修订采纳率和每章总耗时。

\subsection{场景二：存量稿件整理、接续与局部重写}

\textbf{使用情境。}作者有多年积累的 Markdown 或 TXT 稿件，希望在保留既有世界观和人物关系的前提下恢复更新；某次重写又改变了关键事件发生时间，需要核查本次事件对后续章节影响。

\textbf{Agent 协作。}先沿用项目的导入预览与结构确认流程，确认卷章划分并逐步恢复事实；迁移实施中再补齐相应的工具入口。Planner-Agent 整理作者认可的设定，Writer-Agent 从当前有效版本继续写作。局部改稿由 Revision-Agent 结合证据生成修订提案与正文差异，供作者预览，批准后通过领域服务应用；事实接纳流程中，将会标记受影响的章节记忆、人物状态、真相与规划，完成重建后由作者复核。

\textbf{作者获得的结果。}可检查的章节结构、确认后的事实基线、接续候选稿和变更影响清单。候选设定须经作者选择后进入正典；工作流中断时依据操作状态继续处理，保留尚未完成的阶段。

\textbf{价值与验证。}作者可以通过项目检查导入前后正文完整性、章节顺序、被修改事实的影响范围，以及判断写作时是否重复写入。

\subsection{场景三：小型团队的章节审校与定稿协作}

\textbf{使用情境。}作者与编辑共同维护一部作品，需要识别“故事逻辑存在问题”“文字表达可优化”和“当前证据不足”的差别，并确认修改是否确实解决问题。

\textbf{Agent 协作。}三名 Review-Agent 读取同一正文版本，按统一蓝图提交各自负责领域的评分依据与问题。Revision-Agent 围绕选中的问题生成修订提案，作者预览后决定是否采用。应用后旧评审失效，系统以修订后的正文重新核查问题及影响域，再更新后续写作所需上下文。

\textbf{团队获得的结果。}绑定正文版本的评审记录、问题证据、采用或保留的决定、修订前后差异与复评结果，支持编辑交接和后续追查。

\textbf{价值与验证。}以人工标注核查证据定位准确性、误报漏报、复评后问题是否解决，以及是否出现新的回归。
